% =============================================================================
% Section 6 -- Illustrative Case Study: Aria at RetailCo
% =============================================================================
\section{Illustrative Case Study: Aria at RetailCo}
\label{sec:case-study}

We illustrate the APIP framework through a hypothetical scenario grounded in
realistic deployment conditions. RetailCo is a mid-sized e-commerce retailer
that deployed a customer service agent, Aria, to handle order status inquiries,
return requests, and billing disputes. Aria was built on a large language model
with a RAG pipeline indexing RetailCo's product catalog, returns policy, and
billing FAQ documents. At deployment, Aria achieved a task completion rate of
87\% and a CSAT score of 4.3 out of 5.0. All numerical values in this
section---performance thresholds, check-in days, observed metric
movements---are illustrative parameter choices rather than empirical
measurements, selected to be consistent with the simulation in
Section~\ref{sec:simulation}.

\subsection{The Degradation Event}
\label{subsec:degradation-event}

Ninety days post-deployment, RetailCo's customer support team began noticing an
increase in escalation requests from customers who had interacted with Aria. An
informal audit revealed that Aria was consistently citing an outdated version
of the returns policy---specifically, citing a 30-day return window that had
been extended to 45 days in a policy update three weeks prior. Additionally, a
new product line launched six weeks post-deployment was producing a high rate
of ``I don't have information about that product'' responses, reflecting a RAG
index that had not been refreshed.

In the absence of a formal APIP process, the response was: a developer updated
the system prompt to include the new return window figure (a Tier 1
intervention applied to a Tier 2 problem), and a ticket was filed to update the
RAG index. The system prompt change was deployed without a regression test,
inadvertently overriding a constraint that prevented Aria from offering
unsolicited discounts. Over the following week, discount offer rates tripled,
producing a measurable revenue impact before the issue was detected.

\subsection{APIP Application}
\label{subsec:apip-application}

Under the APIP framework, the scenario unfolds differently. RetailCo's
behavioral contract for Aria specifies trigger conditions including: task
completion rate below 82\% (7-day rolling), CSAT below 3.8 (7-day rolling), and
policy citation error rate above 5\% (measured via weekly human review sample of
50 conversations). The policy citation error rate trigger fires at week 10.

Phase 2 cause attribution identifies the failure as behavioral drift: the RAG
index is stale relative to the current policy and product catalog. The failure
is not attributable to model capability, prompt constraints, or tool
configuration. The cause attribution report notes that the system prompt
intervention applied in the ad-hoc scenario was a misclassification of failure
mode.

Phase 3 intervention planning selects Tier 2: RAG index refresh for the returns
policy document and product catalog additions. The intervention plan specifies
a regression test suite to be run before re-deployment, including coverage of
the discount constraint behavior.

A 14-day remediation window is opened with check-ins at days 4, 7, and 11.
Following Kirkpatrick's framework, the day 4 check-in evaluates at multiple
levels: Level 2 (does Aria now cite the correct policy?), Level 3 (has the
pattern of escalation-triggering responses changed?), and Level 4 (is CSAT
recovering?). The policy citation error rate has dropped to below 2\% following
the index refresh. At day 14, the exit evaluation confirms resolution: task
completion has returned to 86\%, CSAT to 4.2, and policy citation errors to
1.5\%. The APIP is closed with outcome: \texttt{resolved}. The full APIP
record---cause attribution report, intervention specification, regression test
results, Kirkpatrick-level check-in assessments, and metric snapshots---is
stored in RetailCo's AI audit trail.

\subsection{Lessons from the Case Study}
\label{subsec:lessons}

The case study illustrates several key properties of the APIP framework. First,
the discipline of cause attribution prevents tier mismatch---applying the
correct intervention tier is as important as applying any intervention. Second,
the behavioral contract formalizes the trigger, making it objective and
consistently applied rather than dependent on informal noticing. Third, the
documented APIP record creates organizational memory: when Aria's RAG index
next requires refreshing, the prior APIP provides a template for the
intervention and a baseline for exit criteria. Fourth, the regression test
requirement embedded in the intervention plan prevents the cascading failure
that occurred in the ad-hoc scenario. Fifth, the Kirkpatrick-level evaluation
at check-ins provides confidence that recovery is durable rather than
superficial.
